\chapter{Operations Research}

\emph{Operations research} is the discipline of applying advanced analytical methods to help
make better decisions. It is used in logistics, finance, healthcare, manufacturing, etc.

\section{Linear Programming}

\emph{Linear programing} si the process od optimizing a linear objective function subsect to a set of linear functional constraints.
Formally, given \(A \in \R^{m \times n}\), \(b \in \R^m\), and \(c \in \R^n\), a linear programming problem can be written as:

\[
    \text{maximize } c^T x \text{ subject to } Ax \le b, \; x \ge 0.
\]

We denote our optimal solution as

\[
    x^{*} = \argmax \{c^T x \mid x \ge 0, Ax \le b\}.
\]

\subsection{Standard Form}

A linear program is in the \emph{standard form} if \(\forall b_i \in b, \ b_i \ge 0\). This means that we can manipulate our
expresions to get \(b\) to the standard form if needed.

We also can modify our problem in the following way

\[
    g(x) \ge b \to -g(x) \le -b
\]

and

\[
    g(x) = b \to g(x) \le b \land g(x) \ge b
\]

This also allows to get an optimal solution from a minimization problem

\[
    \min(f(x)) = \max (-f(x)).
\]

When modelling a problem as a linear program, \textbf{always} use the standard form.

\textbf{Example:}

Given

\[
    z = 3 x_1 + 2 x_2
\]

which we want to maximaze subject to

\begin{align*}
    -x_1 - x_2 &\le -3 \\
    2x_1 + x_2 &\le 8 \\
    -x_1 + x_2 &\le 1 \\
    x_1 + 2x_2 &\le 8 \\
\end{align*}

and the non-negativity constraints \(x_1, x_2 \ge 0\). The \emph{standard form} of this problem is given by

\[
    Ax =
    \begin{bmatrix}
        1 & 1 \\
        2 & 1  \\
        -1 & 1 \\
        1 & 2
    \end{bmatrix}
    \begin{bmatrix}
        x_1 \\
        x_2
    \end{bmatrix}
    \leq
    \begin{bmatrix}
        3 \\
        8 \\
        1 \\
        8
    \end{bmatrix}
    = b.
\]

With

\[
    c =
    \begin{bmatrix}
        3 \\
        2
    \end{bmatrix}.
\]

\section{Vertex Enumeration}

\emph{Vertex enumeration} is the process of listing all vertices (extreme points) of a convex polyhedron defined by a set of linear
inequalities.

Formally, given \(A \in \R^{m \times n}\) and \(b \in \R^m\), we want to find all \(x \in \R^n\) such that:

\[
    Ax \le b.
\]

The vertices of this polyhedron are the points that cannot be expressed as a convex combination of
other points in the polyhedron.

It is computationally challenging and expensive, especially for high-dimensional polyhedra.

\section{Families of Methods for Solving Linear Programs}

\subsection{Simplex Method}

The \emph{simplex method} is an iterative algorithm which ``moves'' along the edges of the feasible region defined
by the linear constraints, from one vertex to another, in search of the optimal solution.

Formally, starting from an initial feasible vertex \(x_0\), the simplex method iteratively moves to an adjacent vertex
\(x_{k+1}\) such that \(c^T x_{k+1} > c^T x_k\), until no further improvement is possible, at which point the current vertex is the optimal solution
\(x^*\).

The simplex method is widely used in practice due to its efficiency on many real-world problems, despite its worst-case
exponential time complexity.

\subsection{Exterior Point Methods}

\emph{Exterior point methods} are a class of algorithms for solving linear and nonlinear convex optimization problems. Unlike the simplex method, which traverses
the boundary of the feasible region, exterior point methods move through the interior and exterior of the feasible region. They are not mainstream.

\subsection{Intetior Point Methods}

\emph{Interior point methods} are a class of algorithms for solving linear and nonlinear convex optimization problems. They operate by traversing the interior of the feasible region,
rather than its boundary, and iteratively improve the solution until convergence to the optimal point.

A common example of an interior point method is the \emph{primal-dual interior point method}, which simultaneously considers the primal and dual formulations of the
optimization problem to guide the search for the optimal solution.

\section{The Simplex Method}

The \emph{simplex method} uses the corner-points of the feasable region to find the best solution. To accomplish this we solve a system of equations induced by a subset
of inequalities of our LP; the \textbf{unique solutions} of this systems are the cornerpoints of our corner points. Also note, that
LPs are system of inequelities, not of equalities threfore, we need to \textbf{modify} our LP get equalities.

We use projections to a higher dimensional space by adding extra \emph{slack variables} so that our inequalities become equalities like for example:

\[
    x_1 \le 4,\,  x_1 \ge 0 \to x_1 + x_3 = 4, \, x_1, x_3 \ge 0.
\]

We do this for all of our functional constraints, giving us the \emph{augmented for of a linear program}.

From this augmented system there are subsets which will have unique solution and hence will give us the corner points.

\subsection{Basic Feasable Solution}

A \emph{basic solution} is a solution which cames from a system of our augmented form which also feasable if and only if all of the decision variables and slack
variables are non-negative.

The variables of our subs-system of equations are called \emph{basic} variables while the remaning are called \emph{non-basic}.

\subsubsection{Adjacency}


\subsubsection{Convexity}


\subsection{The Simplex Algorithm}

Note that we use the value of the objective function

\[
    Z =  a_1 x_1 + \cdots + a_n x_n
\]

as a variable hence,

\[
    Z - a_1 x_1 - \cdots - a_n x_n = 0.
\]

\begin{enumerate}

    \item Find basic feasable solution to start by using the slack variables as a basis for the first system.

    \[
        \begin{bmatrix}
            x_{n + 1} \\
            x_{n + 2} \\
            \vdots \\
            x_{n + m}
        \end{bmatrix}
        =
        \begin{bmatrix}
            b_1 \\
            b_2 \\
            \vdots \\
            b_n
        \end{bmatrix}
    \]

    \item Can a current \textbf{non-basic} variable improve the objective function? Check coefficients in row 0. If no, terminate. 

    \item Pick a non-basic variable with the greatest potential to improbe the objective function as the entering variable.
    st potential to improbe the objective function as the entering variable.
    
    \[
        x_{i} = \max (\{c_i \mid c_i \text{ is coefficient of the target function (also after modification)}\})
    \]
    
    \item Use the minimum ratio test to identify the leaving variable. Using the system of equation gained from the basic solution check which one has the 
    smallest fraction when solving for the \(x_i\). That variable \(x_j\) is the one which is going to hit zero first. This is the \emph{minimum ratio test}

    \begin{align*}
        x_j = \min \left( \frac{b_i}{a_{ij}} \mid a_{ij} > 0 \right)    
    \end{align*}

    \item Update the system for the new basic feasible solution using Gaussian elimination. 
    This means we need to update the column of the variable \(x_i\) for all coefficients to be zero, except for the row \textbf{owned} by 
    the variable the variable \(x_j\) identified by the minimum ratio test as the leaving variable.

    \item Read the solution by noting the solution from the target variables and the slack variables which did not turn 0.

\end{enumerate}

If a our system of equations fails to have a unique solution, then our problem is either infisible or it has infinitely
many solutions. In both cases our simplex is not affected since for the first case we can not start the algorithm and for the second case,
it will just find a solution and then stop.


\textbf{Example:}

Given \(3_x1 + 5 x_2 = Z\) subject to 

\begin{align*}
    x_1  &\le 4 \\ 
    x_2 &\le 12 \\ 
    3x_1 + 2x_2 &\le 18
\end{align*}

We first bring the problem to the standard form 

\begin{align*}
    Z - 3_x1 - 5 x_2 = 0\
    x_1  + x_3 &= 4 \\ 
    x_2  + x_4 &= 12 \\ 
    3x_1 + 2x_2  + x_5 &= 18
\end{align*}


 First, we know that there is a basic solution

\[
    \begin{bmatrix}
        x_3 \\ 
        x_4 \\ 
        x_5
    \end{bmatrix}
    = 
    \begin{bmatrix}
        4 \\ 
        12 \\ 
        18
    \end{bmatrix}
\]

with \(x_1 = 0 = x_2\).

Now, we know that \(x_2\) has the biggest impact on the target function since it has weight \(5\). Therefore, we need 
to take \(x_2\) into consideration for our subsystem of equations and solve

\begin{align*}
    x_3 &= 4 \\ 
    x_4 &= 12 - 2x_2 \\ 
    x_5 &= 18 - 2x_2
\end{align*}

Here since \(\frac{12}{2} < \frac{18}{2}\) we need need to decrease \(x_4\) to zero first.

\[
    \left[
    \begin{array}{ccccc|c}
       -3 & -5 & 0 & 0 & 0 & 0 \\ 
        1 & 0 & 1 & 0 & 0 & 4 \\ 
        0 & 2 & 0 & 1 & 0 & 12 \\   
        3 & 2 & 0 & 0 & 1 & 18 
    \end{array}
    \right]
\]

Row 3 becomes 

\[
    \left[
    \begin{array}{ccccc|c}
       -3 & -5 & 0 & 0 & 0 & 0 \\ 
        1 & 0 & 1 & 0 & 0 & 4 \\ 
        0 & 1 & 0 & \frac{1}{2} & 0 & 6 \\   
        3 & 2 & 0 & 0 & 1 & 18 
    \end{array}
    \right]
\]


Now we need to get ri of the -5 in row 1 and the 2 in  row 4. 

\[
    \left[
    \begin{array}{ccccc|c}
       -3 & 0 & 0 & 2.5 & 0 & 30 \\ 
        1 & 0 & 1 & 0 & 0 & 4 \\ 
        0 & 1 & 0 & \frac{1}{2} & 0 & 6 \\   
        3 & 0 & 0 & -1 & 1 & 6
    \end{array}
    \right]
\]

This completes the first iteration of the algorithm. Note that the non-basic variable \(x_1\) still has potential and that 
the coefficient of \(x_4\) is already non-negative.

We now contine with \(x_1\):

\begin{align*}
        x_3 &= 4 - x_1 \\ 
        x_5 &= 6 - 3x_1
\end{align*}.

Since \(4 > 2\) we know that \(x_5\) has to leave. Hence we will perform Gaussian elimination again in the row of the leaving variable \(x_5\):

Row 3, 

\[
    \left[
    \begin{array}{ccccc|c}
       -3 & 0 & 0 & 2.5 & 0 & 30 \\ 
        1 & 0 & 1 & 0 & 0 & 4 \\ 
        0 & 1 & 0 & \frac{1}{2} & 0 & 6 \\   
        1 & 0 & 0 & \frac{-1}{3} & \frac{1}{3} & 2
    \end{array}
    \right]
\]

For the other rows =: 

\[
    \left[
    \begin{array}{ccccc|c}
        0 & 0 & 0 & \frac{3}{2} & 1 & 36 \\ 
        0 & 0 & 1 & \frac{1}{3} & \frac{-1}{3} & 2 \\ 
        0 & 1 & 0 & \frac{1}{2} & 0 & 6 \\   
        1 & 0 & 0 & \frac{-1}{3} & \frac{1}{3} & 2
    \end{array}
    \right]
\]

Since our slack variables except for \(x_3\) are zero we know that 

\begin{align*}
    Z + 0x_1 + 0x_2 + 0x_3+ \frac{3}{2}x_4 + 1x_5 &= 36 \\
    0x_1 + 0x_2 + x_3 + 0x_4 + 0x_5 &= 2 \\
    0x_1 + x_2 + 0x_3 + 0x_4 + 0x_5 &= 6 \\
    x_1 + 0x_2 + 0x_3 + 0x_4 + 0x_5 &= 2 \\
\end{align*}

which corresponds to 

\begin{align*}
    Z  &= 36 \\
    x_3&= 2 \\
    x_2 &= 6 \\
    x_1 &= 2 \\
    x_4 &= 0 \\
    x_5 &= 0 \\
\end{align*}


Giving us 

\[
    \begin{bmatrix}
        2 \\
        6 \\ 
        2 \\ 
        0 \\ 
        0
    \end{bmatrix}
\]

and the maximum of our objective function 36.
